\chapter*{Реферат}
\addcontentsline{toc}{chapter}{Реферат}
\WorkingTitle\par
\AuthorName\par
\medskip
\noindent\textbf{Ключевые слова:} вызванный составной потенциал действия (ECAP), стимуляция спинного мозга, коннектом, Drosophila melanogaster, нейродинамика, предобучение, перенос представлений.
\medskip
В первой главе рассмотрены наблюдаемые величины Drosophila, анатомические ограничения нейродинамических моделей и механизмы формирования ECAP. Сопоставлены методы моделирования и предобучения, сформулированы требования к исходным данным и независимой проверке переноса. Экспериментальные результаты разрабатываемого метода не представлены.

\ResultPending{Планируется разработать метод биофизического моделирования кривой роста ECAP с предобучением временного модуля на коннектомно-ограниченной динамике Drosophila. Предполагаются отдельные проверки модели нейродинамики и биофизического оператора, а затем оценка дополнительного вклада переноса на независимых записях человека. Область применимости метода будет определена с учётом точности расчётов и неопределённости оценок.}
